% Fichier généré automatiquement par tools/generate_corrections.py.
% Source : SYS/SYS-01/SYS-01-ChainePuissance/1100_Pneumatique/1100_Pneumatique.tex
% Statut : partiel (1/5 question(s) renseignée(s), 0 reprise(s)).
\begin{corrigebox}[Corrigé partiel extrait de la banque]
\CorrigeQuestion{1}
\CorrigeACompleter{Aucun élément de réponse n'est présent dans la source d'origine pour cette question.}
\CorrigeQuestion{2}
\CorrigeACompleter{Aucun élément de réponse n'est présent dans la source d'origine pour cette question.}
\CorrigeQuestion{3}
\CorrigeACompleter{Aucun élément de réponse n'est présent dans la source d'origine pour cette question.}
\CorrigeQuestion{4}
\CorrigeACompleter{Aucun élément de réponse n'est présent dans la source d'origine pour cette question.}
\CorrigeQuestion{5}
Démarrage de la pompe et montée en pression du circuit avec remplissage de l’accumulateur (c).

À la fin de la temporisation le distributeur peut être commandé et ainsi alimenter les vérins.

Si la pression augmente trop, alors le limiteur de pression (d) renvoie une partie du fluide vers le
réservoir et si c’est insuffisant alors (b) permet une décharge du circuit (ouverture vers le réservoir
jusqu’à atteindre ne niveau bas réglé).

La temporisation permet d’attendre qu’un niveau de pression suffisant dans le circuit soit atteint.

Pour remplacer la temporisation on peut mesurer la pression dans le circuit ou plus simplement détecter
le niveau de pression satisfaisant pour le fonctionnement à l’aide d’un pressostat.

La solution utilisant un capteur de pression est plus sûre que la temporisation qui pourrait autoriser la
commande du distributeur alors que la pression dans le circuit est encore insuffisante.

(UPSTI).
\end{corrigebox}
